% GENERATED FILE. Edit canonical lessons, then run npm run generate:books.
% canonical-source-hash: fnv1a64:8ad4049a07c234b4
% canonical-lessons: FR-C317-defile, FR-C317-defiler, FR-C317-carnaval, FR-C317-char, FR-C317-deguiser

\chapter{Le carnaval}
\label{ch:fr-le-carnaval}
\hlchaptermodality{317}

\begin{chapteropening}
\textbf{By the end of this chapter:} \emph{I can describe a carnival: the parade, marching past, the carnival, a float and dressing up.}

Complete the last lesson of chapter 317: I can describe a carnival: the parade, marching past, the carnival, a float and dressing up.
\end{chapteropening}

\section[le défilé]{le défilé — the parade; a fashion show}

\label{lesson:FR-C317-defile}



\begin{quote}
Before the new one: say the French for the brass band, then the French for the public holiday.
\end{quote}

\subsection*{You'll want to know: le défilé}
\textbf{le défilé} — \textquotedblleft{}the parade; a fashion show\textquotedblright{}.

\begin{quote}
Les enfants regardent le défilé du 14 juillet.
\end{quote}

\begin{quote}
\emph{The children are watching the 14 July parade.}
\end{quote}

\emph{Un défilé de mode} — a fashion show.

\begin{cousinweb}
\emph{défilé} is the past participle of \emph{défiler} \textquotedblleft{}to march past\textquotedblright{} used as a noun (EtymDB).
\end{cousinweb}

\begin{grammarlens}[title={What you've built}]
One more word for telling, explaining and describing at B1.
\end{grammarlens}

\subsection*{Your turn}
\begin{itemize}
\raggedright
  \item \emph{Say it:} \emph{le défilé}
  \item \emph{Say it:} \emph{le défilé}, once more
  \item \emph{Say it:} a sentence of your own with \emph{le défilé}
  \item \emph{Recall:} say the French for the brass band, then the French for the public holiday, then say \emph{le défilé} again
\end{itemize}

\begin{tcolorbox}[breakable,colback=teal!4,colframe=teal!35!black,title={Before you move on}]
What does le défilé mean? (\textquotedblleft{}The parade; a fashion show\textquotedblright{}.) Say it once more.
\end{tcolorbox}
\section[défiler]{défiler — to march past, to parade; to scroll}

\label{lesson:FR-C317-defiler}



\begin{quote}
Before the new one: say the French for the public holiday, then the French for the parade.
\end{quote}

\subsection*{You'll want to know: défiler}
\textbf{défiler} — \textquotedblleft{}to march past, to parade; to scroll\textquotedblright{}.

\begin{quote}
Les soldats ont défilé dans la rue.
\end{quote}

\begin{quote}
\emph{The soldiers marched down the street.}
\end{quote}

On a screen: \emph{faire défiler une page} — to scroll.

\begin{cousinweb}
\emph{défiler} is built on \emph{file} \textquotedblleft{}a row, a line\textquotedblright{} (EtymDB), which comes from \emph{fil} \textquotedblleft{}thread\textquotedblright{}, from Latin \emph{fīlum} \textquotedblleft{}a thread\textquotedblright{} (Lewis \& Short).
\end{cousinweb}

\begin{grammarlens}[title={What you've built}]
One more word for telling, explaining and describing at B1.
\end{grammarlens}

\subsection*{Your turn}
\begin{itemize}
\raggedright
  \item \emph{Say it:} \emph{défiler}
  \item \emph{Say it:} \emph{défiler}, once more
  \item \emph{Say it:} a sentence of your own with \emph{défiler}
  \item \emph{Recall:} say the French for the public holiday, then the French for the parade, then say \emph{défiler} again
\end{itemize}

\begin{tcolorbox}[breakable,colback=teal!4,colframe=teal!35!black,title={Before you move on}]
What does défiler mean? (\textquotedblleft{}To march past, to parade; to scroll\textquotedblright{}.) Say it once more.
\end{tcolorbox}
\section[le carnaval]{le carnaval — the carnival}

\label{lesson:FR-C317-carnaval}



\begin{quote}
Before the new one: say the French for the parade, then the French for to march past.
\end{quote}

\subsection*{You'll want to know: le carnaval}
\textbf{le carnaval} — \textquotedblleft{}the carnival\textquotedblright{}.

\begin{quote}
Le carnaval a lieu en février.
\end{quote}

\begin{quote}
\emph{The carnival takes place in February.}
\end{quote}

Plural \emph{les carnavals}.

\begin{cousinweb}
\emph{carnaval} was borrowed from Italian \emph{carnevale} (EtymDB), from Medieval Latin \emph{carnelevāre} \textquotedblleft{}putting away meat\textquotedblright{}, from \emph{carō} \textquotedblleft{}flesh\textquotedblright{} and \emph{levāre} \textquotedblleft{}to lift\textquotedblright{}.
\end{cousinweb}

\begin{grammarlens}[title={What you've built}]
One more word for telling, explaining and describing at B1.
\end{grammarlens}

\subsection*{Your turn}
\begin{itemize}
\raggedright
  \item \emph{Say it:} \emph{le carnaval}
  \item \emph{Say it:} \emph{le carnaval}, once more
  \item \emph{Say it:} a sentence of your own with \emph{le carnaval}
  \item \emph{Recall:} say the French for the parade, then the French for to march past, then say \emph{le carnaval} again
\end{itemize}

\begin{tcolorbox}[breakable,colback=teal!4,colframe=teal!35!black,title={Before you move on}]
What does le carnaval mean? (\textquotedblleft{}The carnival\textquotedblright{}.) Say it once more.
\end{tcolorbox}
\section[le char]{le char — the float (in a parade); the tank}

\label{lesson:FR-C317-char}



\begin{quote}
Before the new one: say the French for to march past, then the French for the carnival.
\end{quote}

\subsection*{You'll want to know: le char}
\textbf{le char} — \textquotedblleft{}the float (in a parade); the tank\textquotedblright{}.

\begin{quote}
Les chars du carnaval sont pleins de fleurs.
\end{quote}

\begin{quote}
\emph{The carnival floats are full of flowers.}
\end{quote}

In Québec, \emph{le char} is also the everyday word for a car.

\begin{cousinweb}
\emph{char} is inherited from Latin \emph{carrus} \textquotedblleft{}a wagon\textquotedblright{} (EtymDB), a word of Gaulish origin.
\end{cousinweb}

\begin{grammarlens}[title={What you've built}]
One more word for telling, explaining and describing at B1.
\end{grammarlens}

\subsection*{Your turn}
\begin{itemize}
\raggedright
  \item \emph{Say it:} \emph{le char}
  \item \emph{Say it:} \emph{le char}, once more
  \item \emph{Say it:} a sentence of your own with \emph{le char}
  \item \emph{Recall:} say the French for to march past, then the French for the carnival, then say \emph{le char} again
\end{itemize}

\begin{tcolorbox}[breakable,colback=teal!4,colframe=teal!35!black,title={Before you move on}]
What does le char mean? (\textquotedblleft{}The float (in a parade); the tank\textquotedblright{}.) Say it once more.
\end{tcolorbox}
\section[déguiser]{déguiser — to dress up (as se déguiser); to disguise}

\label{lesson:FR-C317-deguiser}



\begin{quote}
Before the new one: say the French for the carnival, then the French for the float.
\end{quote}

\subsection*{You'll want to know: déguiser}
\textbf{déguiser} — \textquotedblleft{}to dress up (as se déguiser); to disguise\textquotedblright{}.

\begin{quote}
Les enfants se sont déguisés en animaux.
\end{quote}

\begin{quote}
\emph{The children dressed up as animals.}
\end{quote}

Mostly reflexive in this sense: \emph{se déguiser en} — to dress up as. Without \emph{se}, to disguise: \emph{déguiser sa voix}.

\begin{cousinweb}
\emph{déguiser} comes from Old French \emph{desguiser} (EtymDB), built on \emph{guise} \textquotedblleft{}manner, way\textquotedblright{}, which is usually traced to a Frankish word.
\end{cousinweb}

\begin{grammarlens}[title={What you've built}]
One more word for telling, explaining and describing at B1.
\end{grammarlens}

\subsection*{Your turn}
\begin{itemize}
\raggedright
  \item \emph{Say it:} \emph{déguiser}
  \item \emph{Say it:} \emph{déguiser}, once more
  \item \emph{Say it:} a sentence of your own with \emph{déguiser}
  \item \emph{Recall:} say the French for the carnival, then the French for the float, then say \emph{déguiser} again
\end{itemize}

\begin{tcolorbox}[breakable,colback=teal!4,colframe=teal!35!black,title={Before you move on}]
What does déguiser mean? (\textquotedblleft{}To dress up (as se déguiser); to disguise\textquotedblright{}.) Say it once more.
\end{tcolorbox}
